
Neural network weights encode everything about a model's learned function, yet extracting meaningful properties directly from these high-dimensional parameter vectors remains challenging. Weight-to-accuracy prediction---inferring test performance without running inference---enables efficient model selection, zoo curation, and training diagnostics. Prior work demonstrates this task is solvable: simple per-layer statistics (mean, variance, spectral norm) fed to linear regression achieve $R^2$ > 0.98 on large model collections \citep{unterthiner2020predicting}.

The core difficulty lies in permutation symmetry. Within each layer, hidden units can be arbitrarily reordered (with corresponding permutation of downstream weights) without changing the network's input-output mapping. This symmetry creates exponentially many weight configurations representing identical functions---N! equivalent parameterizations per N-unit layer.

Two architectural strategies address this symmetry. \textbf{Feature engineering} computes permutation-invariant statistics (layer-wise aggregations), discarding structural information but guaranteeing invariance. \textbf{Permutation-equivariant networks} (e.g., Neural Functional Networks) encode symmetry directly in their architecture, processing weights while respecting permutation structure. The equivariant approach preserves richer information but requires specialized layers.

This work investigates a fundamental question: \textbf{Is permutation equivariance merely helpful for data efficiency, or is it categorically required for learning weight-to-accuracy mappings from raw weights?}

Experiments compare three methods across training sizes N $\in$ {100, 250, 500, 1000, 2500, 5000}:
\begin{itemize}
\item \textbf{Statistics baseline}: Per-layer weight statistics with Ridge regression
\item \textbf{MLP baseline}: Flattened raw weights as input features
\item \textbf{NFN model}: Permutation-equivariant DeepSets-style architecture
\end{itemize}

The main finding is that equivariance is not an optimization---it is a prerequisite. The MLP baseline achieves $R^2$ < 0 (worse than mean prediction) at all sample sizes, including N=5000. The equivariant NFN achieves $R^2$ > 0.99 even at N=100. The effect size at N=500 ($\Delta$ = 2.50 $R^2$ points, p < 0.00001) exceeds the predicted threshold by a factor of 25.
